% GENERATED FILE. Edit canonical lessons, then run npm run generate:books.
% canonical-source-hash: fnv1a64:4ce0279235c098c2
% canonical-lessons: FA-C41-miveh, FA-C41-sib, FA-C41-moz, FA-C41-gowje, FA-C41-havij

\chapter{Fruit, Apple, Banana, Tomato, Carrot}
\label{ch:fa-fruit-apple-banana-tomato-carrot}
\hlchaptermodality{41}

\begin{chapteropening}
\textbf{By the end of this chapter:} \emph{I can say fruit, an apple, a banana, a tomato and a carrot.}

Complete the last lesson of chapter 41: I can say fruit, an apple, a banana, a tomato and a carrot.
\end{chapteropening}

\section[\texorpdfstring{\fa{میوه} (miveh)}{میوه (miveh)}]{\fa{میوه} (miveh) — fruit}

\label{lesson:FA-C41-miveh}



\begin{quote}
Before the new one: say the Persian for rice, then the Persian for meat.
\end{quote}

\subsection*{You'll want to know: \fa{میوه}}
\textbf{\fa{میوه}} — \emph{miveh} — \textquotedblleft{}fruit\textquotedblright{}.

\begin{cousinweb}
\emph{miveh} is a native Persian word, inherited rather than borrowed: it continues Middle Persian \emph{mēwag}, \textquotedblleft{}fruit\textquotedblright{} (EtymDB).
\end{cousinweb}

\begin{grammarlens}[title={What you've built}]
One more word for this chapter.
\end{grammarlens}

\subsection*{Your turn}
\begin{itemize}
\raggedright
  \item \emph{Say it:} \emph{miveh}
  \item \emph{Say it:} \emph{miveh}, once more
  \item \emph{Say it:} say \emph{miveh}
  \item \emph{Recall:} say the Persian for rice, then the Persian for meat, then say \emph{miveh} again
\end{itemize}

\begin{tcolorbox}[breakable,colback=teal!4,colframe=teal!35!black,title={Before you move on}]
What does \fa{میوه} mean? (\textquotedblleft{}Fruit\textquotedblright{}.) Say it once more.
\end{tcolorbox}
\section[\texorpdfstring{\fa{سیب} (sib)}{سیب (sib)}]{\fa{سیب} (sib) — an apple}

\label{lesson:FA-C41-sib}



\begin{quote}
Before the new one: say the Persian for meat, then the Persian for fruit.
\end{quote}

\subsection*{You'll want to know: \fa{سیب}}
\textbf{\fa{سیب}} — \emph{sib} — \textquotedblleft{}an apple\textquotedblright{}.

\begin{cousinweb}
\emph{sib} is a native Persian word, inherited rather than borrowed; it is usually traced to Middle Persian \emph{sēb}, \textquotedblleft{}apple\textquotedblright{}.
\end{cousinweb}

\begin{grammarlens}[title={What you've built}]
One more word for this chapter.
\end{grammarlens}

\subsection*{Your turn}
\begin{itemize}
\raggedright
  \item \emph{Say it:} \emph{sib}
  \item \emph{Say it:} \emph{sib}, once more
  \item \emph{Say it:} say \emph{sib}
  \item \emph{Recall:} say the Persian for meat, then the Persian for fruit, then say \emph{sib} again
\end{itemize}

\begin{tcolorbox}[breakable,colback=teal!4,colframe=teal!35!black,title={Before you move on}]
What does \fa{سیب} mean? (\textquotedblleft{}An apple\textquotedblright{}.) Say it once more.
\end{tcolorbox}
\section[\texorpdfstring{\fa{موز} (moz)}{موز (moz)}]{\fa{موز} (moz) — a banana}

\label{lesson:FA-C41-moz}



\begin{quote}
Before the new one: say the Persian for fruit, then the Persian for an apple.
\end{quote}

\subsection*{You'll want to know: \fa{موز}}
\textbf{\fa{موز}} — \emph{moz} — \textquotedblleft{}a banana\textquotedblright{}.

\begin{cousinweb}
\emph{moz} continues Middle Persian \emph{mōz}, \textquotedblleft{}banana\textquotedblright{} (EtymDB). It is usually linked with Sanskrit \emph{mōca}, the plantain and its fruit (CDIAL 10348; EtymDB lists the related \emph{mauca}). Arabic \emph{mawz} is the same word.
\end{cousinweb}

\begin{grammarlens}[title={What you've built}]
One more word for this chapter.
\end{grammarlens}

\subsection*{Your turn}
\begin{itemize}
\raggedright
  \item \emph{Say it:} \emph{moz}
  \item \emph{Say it:} \emph{moz}, once more
  \item \emph{Say it:} say \emph{moz}
  \item \emph{Recall:} say the Persian for fruit, then the Persian for an apple, then say \emph{moz} again
\end{itemize}

\begin{tcolorbox}[breakable,colback=teal!4,colframe=teal!35!black,title={Before you move on}]
What does \fa{موز} mean? (\textquotedblleft{}A banana\textquotedblright{}.) Say it once more.
\end{tcolorbox}
\section[\texorpdfstring{\fa{گوجه} (gowje)}{گوجه (gowje)}]{\fa{گوجه} (gowje) — a tomato}

\label{lesson:FA-C41-gowje}



\begin{quote}
Before the new one: say the Persian for an apple, then the Persian for a banana.
\end{quote}

\subsection*{You'll want to know: \fa{گوجه}}
\textbf{\fa{گوجه}} — \emph{gowje} — \textquotedblleft{}a tomato\textquotedblright{}.

\begin{cousinweb}
\emph{gowje} is short for \emph{gowje farangi}, \textquotedblleft{}European plum\textquotedblright{}: the tomato was named after the plum, \emph{gowje}, and labelled \emph{farangi}, \textquotedblleft{}Frankish, European\textquotedblright{}. Where \emph{gowje} itself comes from is not settled.
\end{cousinweb}

\begin{grammarlens}[title={What you've built}]
One more word for this chapter.
\end{grammarlens}

\subsection*{Your turn}
\begin{itemize}
\raggedright
  \item \emph{Say it:} \emph{gowje}
  \item \emph{Say it:} \emph{gowje}, once more
  \item \emph{Say it:} say \emph{gowje}
  \item \emph{Recall:} say the Persian for an apple, then the Persian for a banana, then say \emph{gowje} again
\end{itemize}

\begin{tcolorbox}[breakable,colback=teal!4,colframe=teal!35!black,title={Before you move on}]
What does \fa{گوجه} mean? (\textquotedblleft{}A tomato\textquotedblright{}.) Say it once more.
\end{tcolorbox}
\section[\texorpdfstring{\fa{هویج} (havij)}{هویج (havij)}]{\fa{هویج} (havij) — a carrot}

\label{lesson:FA-C41-havij}



\begin{quote}
Before the new one: say the Persian for a banana, then the Persian for a tomato.
\end{quote}

\subsection*{You'll want to know: \fa{هویج}}
\textbf{\fa{هویج}} — \emph{havij} — \textquotedblleft{}a carrot\textquotedblright{}.

\begin{cousinweb}
\emph{havij} is treated as a native Persian word, but no older form is given in the sources used here, and its origin is not settled.
\end{cousinweb}

\begin{grammarlens}[title={What you've built}]
One more word for this chapter.
\end{grammarlens}

\subsection*{Your turn}
\begin{itemize}
\raggedright
  \item \emph{Say it:} \emph{havij}
  \item \emph{Say it:} \emph{havij}, once more
  \item \emph{Say it:} say \emph{havij}
  \item \emph{Recall:} say the Persian for a banana, then the Persian for a tomato, then say \emph{havij} again
\end{itemize}

\begin{tcolorbox}[breakable,colback=teal!4,colframe=teal!35!black,title={Before you move on}]
What does \fa{هویج} mean? (\textquotedblleft{}A carrot\textquotedblright{}.) Say it once more.
\end{tcolorbox}
